% M10, 11.09.2026: Leserführung, Kapiteltrennung und belegte Reichweite.
% E47-01: Anhang A, Originale 20260727_074919_b2c654 / 20260727_080638_9d4ae8.
% A9/A10 wortgleich erhalten; Herkunft technische bzw. betriebliche Konsequenz.

\section{Durabilität und kontinuierliche Nutzung}
\label{sec:durabilitaet-kontinuierliche-nutzung}

Die letzten beiden Anforderungen betreffen den Betrieb des gewachsenen
Systems. A9 folgt aus den menschlichen Entscheidungspunkten und den
extern wirksamen Verarbeitungsschritten; A10 entstand aus der
Bedienerfahrung während der Entwicklung.

\emph{Technische Konsequenz.} Die zuvor eingeführten
Autorisierungspunkte machen das Warten auf eine menschliche Entscheidung
zu einem vorgesehenen Ablaufzustand. Diese Wartezeit kann die Lebensdauer
des ausführenden Prozesses überschreiten. Deshalb muss über dessen
Arbeitsspeicher hinaus festgehalten werden, wo ein Lauf steht und welche
Antwort noch aussteht. Dieser Ausführungszustand beantwortet die
Wiederaufnahmefrage; der Core hält dagegen fest, welche fachlichen
Informationen im Projekt gelten
(Abschnitte~\ref{sec:zustand-persistenz-projektwissen}
und~\ref{sec:projektzustand-kontrollierte-mutation}).

Die Gestaltung setzte dafür auf Sicherungspunkte und explizite
Wiederaufnahme. Zwei dokumentierte Machbarkeitsfälle zeigen das Anhalten
an einem menschlichen Entscheidungspunkt und die Fortsetzung nach
Übergabe der Antwort: einmal bis zum Abschluss, einmal bis zum nächsten
vorgesehenen Entscheidungspunkt (E47-01). Daraus folgt noch kein
Nachweis des Schutzes vor Doppelwirkungen. Dieser ist eine zusätzliche
technische Anforderung: Bereits angelegte externe Arbeitselemente sollen
bei einer Wiederaufnahme nicht unkontrolliert erneut entstehen.

\textbf{A9:} Langlaufende und an menschlichen Entscheidungspunkten
unterbrochene Prozesse müssen pausierbar, wiederaufnehmbar und gegen
unbeabsichtigte Doppelwirkung abgesichert sein.

\emph{Betriebliche Konsequenz.} Mit mehreren Verarbeitungsabläufen,
Eingangswegen und Autorisierungspunkten wurde die direkte Bedienung
aufwendiger. Der Bediener musste den passenden Ablauf auswählen, wartende
Läufe finden und ausstehende Entscheidungen zuordnen. Dies ist eine
Beobachtung aus der Entwicklungs- und Stabilisierungsphase mit dem Autor
als einzigem Bediener, keine Erkenntnis aus einer Nutzerstudie.

Als Antwort wurde eine dialogorientierte Steuerungsschicht ergänzt:
Sie interpretiert ein Nutzeranliegen, berücksichtigt den verfügbaren
Systemzustand und wählt aus den bereitgestellten Werkzeugen eine passende
Handlung. Darin liegt ihr agentischer Beitrag. Die aufgerufenen Abläufe
führen weiterhin die fachliche Verarbeitung und ihre Kontrollen aus;
die Bedienschicht implementiert diese Logik nicht erneut. Sie ergänzt
die weiterhin direkt ausführbaren Workflows um einen dialogischen Zugang
und wird in Abschnitt~\ref{sec:durable-ausfuehrung-agentische-bedienung}
als \emph{Steward} konkretisiert. Für fachliche Änderungen nutzt sie die
vorgesehenen Prüf-, Entscheidungs- und Speicheroperationen.

\textbf{A10:} Die kontinuierliche Nutzung des entwickelten Systems soll
durch eine agentische Bedien- und Orchestrierungsschicht unterstützt
werden, die bestehende Governance nutzt, aber nicht ersetzt.

\emph{Einordnung.} Die technische Realisierung folgt in
Kapitel~\ref{chap:maf-realisierung}; die abgegrenzten Nachweise stehen
in den Abschnitten~\ref{sec:eval-kontrollen} und~\ref{sec:eval-resume}.
Sie untersuchen ausgewählte Wiederaufnahme- und Wiederholungsfälle,
keine beliebige Ausfalltoleranz oder doppelte externe Schreibwirkungen
nach einem Absturz. Die Steward-Aufgaben dienen der Demonstration;
eine Usability- oder Produktivitätsbewertung liegt nicht vor. Der
folgende Abschnitt konsolidiert die zehn Designanforderungen und
übergibt sie an die Lösungskonzeption.
